
\section{Experiments and Results}
\label{sec:results}

We organize our experiments around four questions:
(i)~can processing cues alone explain strong performance on
morphing benchmarks?
(ii)~do detectors distinguish identity mixing from other
manipulation artifacts?
(iii)~can training with matched controls improve this
discrimination?
and (iv)~do the resulting improvements transfer to
conventional benchmarks and unseen image sources?

We distinguish three types of evidence throughout:
conventional morph-versus-photograph discrimination,
responses to matched manipulation controls, and
controlled training comparisons. Unless otherwise
specified, results concern individual networks rather
than ensembles. Development-exposed analyses are
identified explicitly and are not treated as independent
confirmation.x

\subsection{Can Morph Benchmarks Be Solved Without Morphs?}
\label{sec:artifact_results}

\paragraph{Morph-free training.}
We first investigate whether a morphing benchmark can
be solved without learning from any morph examples.
To this end, we train networks to distinguish synthetic
SMDD faces~\cite{damer2022smdd} from copies subjected
to downsampling and subsequent enlargement.
Neither training class contains morphed images.

The experiment comprises two CLIP-based architectures,
ViT-B/16 and ViT-L/14, each trained with five independent
seeds. Each training set contains 25,000 original synthetic
faces and 25,000 corresponding processed copies.
The processing pipeline reduces images from $256$ to
$80$ pixels using area interpolation, stores them as
lossless PNG files, and subsequently resizes them to
$224$ pixels using bilinear interpolation.
Both architectures are trained for 6,250 updates,
and final checkpoints are evaluated without
validation-based selection.

All ten networks achieve
\textbf{0.0\% EER} on the native MorDIFF evaluation
under our preprocessing~\cite{damer2023mordiff},
despite never observing a morph during training.
This demonstrates that processing-related differences
alone are sufficient for perfect discrimination on this
specific benchmark evaluation. It does not imply that
every morph-trained detector uses the same cue, nor
does it identify blur as the unique explanation:
downsampling and enlargement jointly alter multiple
image properties.

\paragraph{Sensitivity of released detectors.}
We next examine whether established morph detectors
respond to differences in the processing history of
genuine photographs. We evaluate the released
checkpoints of SelfMAD~\cite{ivanovska2025selfmad},
MADation~\cite{caldeira2025madation}, and
MixFaceNet-MAD using the controls introduced in
Section~\ref{sec:matched_controls}.

On corrected MAD22, replacing raw genuine photographs
with aligned-history controls changes SelfMAD's
WebMorph EER from $13.73\%$ to $29.90\%$.
Its four-family mean EER increases from $4.72\%$
to $9.89\%$. MADation ViT-B/16 and MixFaceNet-MAD
also exhibit substantial changes, although the
magnitude and direction depend on the model and
morph family.
These interventions modify alignment, crop or field
of view, and processing history together. Therefore,
they demonstrate sensitivity to the combined
intervention rather than isolating a single resampling
or compression artifact.
Together with morph-free benchmark saturation,
these results show that successful conventional
morph-versus-photograph discrimination does not,
by itself, establish morph-specific detection.

% FIGURE PLACEMENT:
% The existing processing/shortcut verdict-map figure
% can be introduced here and placed in Section 5.2,
% where its complete hierarchy is discussed.
% Retain the original graphic and label when moving it.


\subsection{Do Detectors Recognize Blending or Identity Mixing?}
\label{sec:blend_results}

The previous experiments establish that non-identity
cues can support morph detection. We now use the
matched controls from Section~\ref{sec:controlled_eval}
to investigate whether detector responses persist
when rendering, warping, blending, or contributor
identity are controlled more closely.

\paragraph{Generative rendering versus identity mixing.}
GAN-generated morphs introduce a particular ambiguity:
a detector may recognize artifacts of the generative
pipeline rather than evidence that two identities
contributed to the image. We examine this distinction using MIPGAN-style
constructions comprising 227 camera photographs,
238 single-photo renders, 119 same-person morphs,
and 74 two-person morphs after excluding
training-linked partners. For the five networks trained with our full protocol
(Sb), discrimination between single-photo renders
and camera photographs reaches AUC values of
$0.992$--$0.998$.
Moreover, $78$--$90\%$ of single-photo renders
are classified as morphs at the detector's own
decision threshold, despite containing no
second identity.

In contrast, the AUC for distinguishing two-person
from same-person morphs ranges from only
$0.619$ to $0.690$.
Thus, the detectors respond strongly to the
generative rendering process itself, while
contributor-dependent discrimination is
considerably weaker. The latter remains above
chance in these comparisons, so we cannot conclude
that identity-related information is entirely absent.
However, the observed differences do not establish
that the networks recognize identity mixing rather
than residual generation or low-level statistical
differences.

\paragraph{Blending versus geometric warping.}
We next examine whether detectors identify information
introduced by blending beyond the geometric
transformations required for morph construction. We construct matched FaceMorpher blends and warp-only
controls using 299 FFHQ development hosts.
To reduce dependence on the surrounding image,
we retain the face region and the morphing tool's
halo while replacing the remaining image with gray.
This diagnostic was registered after the original
candidate's benchmark evaluation and is therefore
development-exposed.% All three training configurations discriminate
% blended morphs from warp-only controls
% (Table~\ref{tab:blend_controls}).
The full protocol reaches an AUC of $0.889$,
compared with $0.879$ for processing-matched
training (B) and $0.864$ for the plain baseline (A).
%
However, the improvement of Sb over B is only
$+0.011$ AUC
%, with a 95\% interval of
$[-0.010,0.031]$. The corresponding paired
seed-level interval also includes zero.
Hence, the additional benefit is not established.

% \begin{table}[t]
% \centering
% \small
% \caption{
% Blend-versus-warp discrimination on the face and
% morphing halo. Intervals use host-block bootstrap
% resampling over 299 development hosts and condition
% on the fitted networks.
% }
% \label{tab:blend_controls}
% \begin{tabular}{lcc}
% \toprule
% Training & AUC & 95\% interval \\
% \midrule
% Plain (A)             & 0.864 & [0.848, 0.881] \\
% Protocol controls (B) & 0.879 & [0.862, 0.894] \\
% Full protocol (Sb)    & 0.889 & [0.873, 0.908] \\
% \midrule
% Sb $-$ B & +0.011 & [$-$0.010, +0.031] \\
% \bottomrule
% \end{tabular}
% \end{table}



    \begin{table*}[t]
    \centering\small
    \setlength{\tabcolsep}{2.4pt}
    \caption{MAD22 under explicit processing controls, EER (\%). 
    % Each arrow compares the same checkpoint(s), not retraining. 
    Panel A changes genuine-photo processing; Panel B additionally equalises the border for Sb. These test sensitivity to specified processing mismatches and do not remove all shortcuts. 
    RW = released weights re-evaluated by us, and Sb averages five seeds 403--407.
    % All five morph families are retained.
    % LM2 averages OpenCV/Webmorph; M4 adds MIP-I/II.
    % FaceMorpher is shortcut-saturated, and GAN discrimination is not established morph-specific.
    % MAD22 informed development; its genuine side exposes FRLL photographs.
    Panel C retrains each comparator with our matched-control recipe (five seeds 401--405; pre-registered protocol AD6); its cells average the five seeds.
    Bold marks the best history-matched (H) EER per family, underline -- the second best, with lower better: Panel A is compared on the H value of each arrow, Panel C on its averaged EER, all eight method rows together.
    }
    \label{tab:mad22_results}
    \begin{tabular}{lrrrrrrr}
    \toprule
    Method / source & FM & OpenCV & Webmorph & MIP-I & MIP-II & LM2 & M4 \\
    \midrule
    \multicolumn{8}{l}{\textbf{A. Raw genuine $\rightarrow$ aligned-history genuine (S $\rightarrow$ H)}} \\
\textbf{    Sb (Ours) (5 seeds)} & $2.32\!\to\!3.67$ & $2.28\!\to\!\textbf{2.26}$ & $10.69\!\to\!\textbf{10.39}$ & $0.57\!\to\!\textbf{0.57}$ & $0.28\!\to\!\textbf{0.28}$ & $6.48\!\to\!\textbf{6.32}$ & $3.45\!\to\!\textbf{3.37}$ \\
    MADation B/16 (RW) & $0.00\!\to\!13.60$ & $3.96\!\to\!6.61$ & $11.00\!\to\!14.22$ & $40.20\!\to\!40.70$ & $28.43\!\to\!28.33$ & $7.48\!\to\!10.41$ & $20.90\!\to\!22.46$ \\
    MADation L/14 (RW) & $0.49\!\to\!\underline{2.45}$ & $4.78\!\to\!4.57$ & $28.43\!\to\!24.60$ & $32.30\!\to\!30.00$ & $16.32\!\to\!14.91$ & $16.60\!\to\!14.59$ & $20.46\!\to\!18.52$ \\
    MixFaceNet-MAD (RW) & $4.90\!\to\!13.40$ & $8.43\!\to\!11.99$ & $18.20\!\to\!24.20$ & $16.30\!\to\!21.57$ & $20.59\!\to\!25.98$ & $13.32\!\to\!18.10$ & $15.88\!\to\!20.94$ \\
    SelfMAD (RW) & $0.00\!\to\!\textbf{0.98}$ & $1.02\!\to\!\underline{2.94}$ & $13.73\!\to\!29.90$ & $2.45\!\to\!\underline{4.00}$ & $1.70\!\to\!\underline{2.70}$ & $7.37\!\to\!16.42$ & $4.72\!\to\!\underline{9.89}$ \\
    \midrule
    \multicolumn{8}{l}{\textbf{B. History-matched $\rightarrow$ border-equalised (H $\rightarrow$ BE\_H)}} \\
   \textbf{ Sb (Ours) (5 seeds)} & $3.67\!\to\!4.46$ & $2.26\!\to\!2.26$ & $10.39\!\to\!10.39$ & $0.57\!\to\!0.95$ & $0.28\!\to\!0.50$ & $6.32\!\to\!6.32$ & $3.37\!\to\!3.52$ \\
    \midrule
    \multicolumn{8}{l}{\textbf{C. Comparators retrained with our matched-control recipe (Sb), 5 seeds (401--405)}} \\
    MADation B/16 & $1.89$ & $4.84$ & $\underline{11.48}$ & $17.17$ & $6.48$ & $\underline{8.16}$ & $9.99$ \\
    MADation L/14 & $6.25$ & $4.68$ & $12.12$ & $18.43$ & $11.50$ & $8.40$ & $11.68$ \\
    MixFaceNet-MAD & $6.33$ & $16.68$ & $22.84$ & $29.33$ & $29.00$ & $19.76$ & $24.46$ \\
    \bottomrule
    \end{tabular}
    \par\smallskip\raggedright
    % Morph counts FM/OpenCV/Webmorph/MIP-I/MIP-II: 1000/984/500/1000/999, plus 204 genuine photos of 102 persons. H combines alignment, crop/field of view and processing history, so it does not isolate blur or JPEG. 
    
    % Under the registered release crops, MADation B/L and MixFaceNet BE rows equal their S/H rows and are not new blanked-image evaluations. 
    
    % SelfMAD has no admitted BE result and is not a comparator in Panel B. 
    
    % Full separate rows are retained in the supplement, and per-family FRLL-Morphs cells were afterwards read separately under the 15 November 2026 author ruling (Table~\ref{tab:frll_results}). 
    
    
    % Arrow first values and Panel B are unmarked. Panel C rows are pre-registered matched-control-recipe retrains of the comparators (prereg AD6), five seeds 401-405, scored on the frozen release crops; per-seed EERs and receipts: smad-next experiments/AGENTIC_delivery_20261114/AD6_H_SBR_EER_20261117/. Bold/underline re-evaluated over the enlarged method set (H6).
    \end{table*}
    

To examine alternative explanations, we additionally
evaluate simple image-statistic references.
Cross-fitted statistical models achieve AUC values
of $0.925$ and $0.917$ on these evaluation images.
These are in-domain diagnostic references rather than
directly comparable out-of-domain baselines.
Nevertheless, they demonstrate that local image
statistics remain a viable explanation for
blend-versus-warp discrimination.
Furthermore, masking the surrounding image itself
changes detector behavior: $47$--$99\%$ of
warp-only images are flagged under the face-region
diagnostic, compared with $2$--$14\%$ in their
whole-frame form.
We therefore interpret this experiment as evidence
of sensitivity to blending-related image properties,
not as confirmation of identity-specific detection
or fixed-threshold robustness.

\paragraph{What do the controls establish?}
Across processing, rendering, and blending comparisons,
we find that high morph scores can arise without
combining two distinct identities. The controls progressively constrain alternative
explanations, but none individually identifies a
unique recognition mechanism. In particular,
same-person and two-person morphs can still differ
in texture, noise averaging, geometric displacement,
or manipulation strength.

% INSERT EXISTING VERDICT-MAP FIGURE HERE.
%
% Suggested caption:
% "Summary of the matched-control evaluation.
% Processing, rendering, blending, splice construction,
% and contributor identity represent progressively more
% restrictive diagnostic questions. The results identify
% several non-identity cues used by detectors, while
% identity-specific detection remains unestablished.
% Each verdict applies only to its registered control
% and evaluation population."
%
% Use the original figure source and preserve its label.


\subsection{Does Matched-Control Training Improve Discrimination?}
\label{sec:training_results}

We now evaluate whether the controls that expose
shortcut dependencies can also improve detector
training, following the protocol introduced in
Section~\ref{sec:matched_training}.

%All primary comparisons use CLIP ViT-B/16 with LoRA, a fixed optimization procedure, and five independent seeds per configuration (408--412). We compare the plain baseline (A), processing-control training(B),  splice-positive training (Sp), and the full matched-control protocol (Sb), together with the count-balanced variants defined in Section~\ref{sec:training_comparisons}.

\paragraph{Do matched negatives improve contributor
discrimination?}
Our primary comparison evaluates the full protocol
(Sb) against SpV, which replaces matched splice
negatives with repeated ordinary negatives.
Both configurations contain 8,386 balanced
training rows and identical positive-example
multisets.
We measure discrimination between two-person and
same-person whole-frame morphs using the
$D$-stratified evaluation defined in
Section~\ref{sec:control_protocol}.
This analysis combines 100 development hosts with
80 screened test hosts and accounts for shared
host--partner components during bootstrap
resampling. 

Sb improves two-person versus same-person
whole-frame ranking by
\textbf{+0.0481 AUC}
%, with a 95\% interval of $[0.0269,0.0722]$.
All five paired seed differences are positive. For face-spliced morphs, the corresponding
unstratified improvement is
\textbf{+0.0818 AUC}
%, with a 95\% interval of $[0.0535,0.1083]$.
The comparison without splice positives (BN versus BV)
does not show the same improvement. Its whole-frame
effect is $-0.0030$ AUC, with an interval crossing
zero, while its face-splice effect is negative.
This pattern suggests that matched negatives may
be particularly useful when paired with corresponding
two-person splice positives.

However, the comparisons do not isolate a universal
interaction between these components: the positive
training distributions match within each pair but
not across all variants.
Furthermore, although class balance, training-row
counts, and positive-example exposure are controlled,
Sb contains 312 additional unique images relative
to SpV. The improvement can therefore be attributed
to the combined intervention, but not exclusively
to the semantic content of matched negatives.

\paragraph{Same-person false alarms.}
Improved ranking does not necessarily translate
into correct classification at a fixed threshold.
For face-spliced same-person controls, the fraction
classified as attacks decreases from $70$--$74\%$
with Sp to $36$--$41\%$ with Sb.
For whole-frame same-person morphs, however,
Sb still flags approximately $85\%$ at its own
decision threshold, compared with $92\%$ for B
and $98\%$ for A.
Thus, matched-control training improves relative
discrimination but does not establish reliable
rejection of manipulations involving only one identity.

\paragraph{Storage sensitivity.}
We also investigate whether the proposed training
examples reduce responses to image-storage history.
Under paired processing perturbations at the
trained small-image size, the plain baseline
exhibits genuine false-alarm increases of
$26.5$--$38.5$ percentage points, compared
with $0$--$4.5$ points for Sb.
At an untrained image size, the corresponding
ranges are $5.5$--$15.2$ points for A and
$0$--$2.1$ points for Sb.
These findings demonstrate substantial attenuation
of processing-induced false alarms.
Nevertheless, the registered robustness criterion
also requires acceptable two-sided score movement.
Sb satisfies the criterion in five of five seeds
at the trained size, but only three of five at
the untrained size. The combined verdict is
therefore inconclusive.
We consequently claim reduced sensitivity to
the tested storage perturbations, rather than
complete removal of processing shortcuts.

\paragraph{Component ablations.}
Additional experiments remove warp-only negatives,
individual warped halves, and stored-small copies
from the full training protocol. Removing stored-small copies increases the measured
trained-size false-alarm effect by $0.054$%, with a 95\% interval of $[0.038,0.070]$.
However, adding stored-small copies alone does
not satisfy the registered robustness criterion.
The necessity of warp-only negatives and native
warped halves is not established under the
registered margins.
We report the full component-ablation table,
per-seed outcomes, and additional unsuccessful
training objectives in the supplementary material.
These include experiments with tail-risk objectives,
paired ranking losses, matched hard negatives,
and alternative training-data compositions.


\subsection{Generalization to Existing Benchmarks and
Unseen Sources}
\label{sec:benchmark_results}

We finally examine whether the proposed training
protocol retains conventional morph-detection
performance and whether its operating threshold
transfers to previously unseen genuine-image sources. These experiments complement the matched-control
evaluations but answer a different question:
conventional benchmark discrimination does not
establish identity-specific detection.

\paragraph{Corrected MAD22 evaluation.}
We evaluate five morph families from corrected
MAD22: FaceMorpher, OpenCV, WebMorph, MIPGAN-I,
and MIPGAN-II. The evaluation includes 204 genuine FRLL photographs
of 102 individuals and 4,483 morphed images.
We compare our frozen Sb networks with released
MADation~\cite{caldeira2025madation},
SelfMAD~\cite{ivanovska2025selfmad}, and
MixFaceNet-MAD checkpoints. Table~\ref{tab:mad22_results} reports EER under
raw-genuine and aligned-history genuine controls.
The intervention changes alignment,
field of view, and processing history.
Panel~C of Table~\ref{tab:mad22_results} retrains the
three comparators with our matched-control
recipe (five seeds per comparator) under the
pre-registered protocol. Retraining recovers part of
the comparator gap for MADation: B/16 improves on its
released history-matched EER on all five families
(e.g. $17.17$ versus $40.70$ on MIPGAN-I), and L/14
improves on three of the five families
(WebMorph, MIPGAN-I, MIPGAN-II). The
pattern does not transfer to MixFaceNet-MAD, whose
retrained EER is higher than its released checkpoint
on OpenCV, MIPGAN-I, and MIPGAN-II -- its released
advantage does not survive retraining under a
history-matched recipe. The ranking does not invert:
our frozen Sb networks remain best on six of the seven
columns, and the second-best method changes only on
WebMorph and LM2, where retrained MADation B/16
narrows the margin without overtaking. The
GAN-generated families remain hardest for every
retrained MADation and MixFaceNet row, consistent
with the shortcut analysis above.
Figure~\ref{fig:families_photos} shows one
identity-matched source--morph--source row for
each family.

% \begin{table*}[t]
% \centering
% \small
% \setlength{\tabcolsep}{3.5pt}
% \caption{
% MAD22 EER (\%) corrected under explicit genuine-image
% processing controls. Each entry shows raw genuine
% $\rightarrow$ aligned-history genuine evaluation
% for the same checkpoint(s). Sb reports the mean
% over five seeds; RW denotes released checkpoints.
% LM2 averages OpenCV and WebMorph, while M4 additionally
% includes MIPGAN-I/II. FaceMorpher is excluded from
% these aggregates because processing-only training
% saturates that benchmark cell.
% Results do not establish
% identity-specific detection.
% }
% \label{tab:mad22_results}
% \begin{tabular}{lrrrrrrr}
% \toprule
% Method & FM & OpenCV & WebMorph &
% MIP-I & MIP-II & LM2 & M4 \\
% \midrule
% Sb (ours) (5 seeds)
% & $2.32\to3.67$
% & $2.28\to2.26$
% & $10.69\to10.39$
% & $0.57\to0.57$
% & $0.28\to0.28$
% & $6.48\to6.32$
% & $3.45\to3.37$ \\

% MADation B/16 (RW)
% & $0.00\to13.60$
% & $3.96\to6.61$
% & $11.00\to14.22$
% & $40.20\to40.70$
% & $28.43\to28.33$
% & $7.48\to10.41$
% & $20.90\to22.46$ \\

% MADation L/14 (RW)
% & $0.49\to2.45$
% & $4.78\to4.57$
% & $28.43\to24.60$
% & $32.30\to30.00$
% & $16.32\to14.91$
% & $16.60\to14.59$
% & $20.46\to18.52$ \\

% MixFaceNet-MAD (RW)
% & $4.90\to13.40$
% & $8.43\to11.99$
% & $18.20\to24.20$
% & $16.30\to21.57$
% & $20.59\to25.98$
% & $13.32\to18.10$
% & $15.88\to20.94$ \\

% SelfMAD (RW)
% & $0.00\to0.98$
% & $1.02\to2.94$
% & $13.73\to29.90$
% & $2.45\to4.00$
% & $1.70\to2.70$
% & $7.37\to16.42$
% & $4.72\to9.89$ \\
% \bottomrule
% \end{tabular}
% \end{table*}

Sb exhibits comparatively small changes under
this particular processing intervention.
Its four-family mean changes from $3.45\%$
to $3.37\%$, whereas SelfMAD's mean changes
from $4.72\%$ to $9.89\%$.
However, the difference between Sb and the strongest
released comparator is $-1.4$ EER percentage
points%, with a two-level 95\% interval of
$[-3.51,0.98]$.
Superiority is therefore not established.
%
Moreover, the proposed recipe was selected after
earlier MAD22 results had been inspected.
Although the reported Sb checkpoints were frozen
before their own benchmark read, this remains a
development-exposed comparison, not an untouched
confirmation of generalization.
%
The low MIPGAN errors also cannot be interpreted
as evidence of identity-specific detection in
light of the rendering-control results of
Section~\ref{sec:blend_results}.

\paragraph{FRLL-Morphs.}
We additionally evaluate the frozen Sb networks
and released comparators on five FRLL-Morphs
families: aMSL, FaceMorpher, OpenCV, StyleGAN,
and WebMorph.
For Sb, the mean EERs are respectively
$11.85\%$, $2.35\%$, $1.37\%$, $10.87\%$,
and $11.69\%$.
The results are mixed across families.
For example, SelfMAD achieves $0.16\%$ EER
on OpenCV, compared with $1.37\%$ for Sb,
while MixFaceNet-MAD achieves $9.08\%$
on StyleGAN, compared with $10.87\%$ for Sb.
This evaluation is descriptive: its genuine
photographs overlap those used in MAD22,
and the corresponding identities were already
exposed during model development.
We therefore do not claim independent generalization
or state-of-the-art performance from these results.

The complete per-family table, released-checkpoint
comparisons, and original-recipe retraining results
are provided in the supplementary material.

\paragraph{Fixed-threshold transfer to unseen sources.}
EER measures discrimination with an evaluation-dependent
operating point. In deployment, however, a detector
must classify images from previously unseen sources
without adjusting its decision threshold using
their labels.
%
We therefore examine genuine-image false alarms
at each detector's fixed decision rule using
previously unseen photographic sources.
%
On eight web-image sources, the behavior varies
substantially across training configurations.
Some configurations exhibit substantial shifts
in score distributions even when within-source
ranking remains strong.

We additionally evaluate 865 frontal NIST MEDS
mugshots from 400 individuals.
False alarms concentrate in a subgroup of
110 small $240\times240$ photographs whose
faces are enlarged by approximately three times.
%
On this subgroup, the plain baseline flags up
to $39\%$ of genuine photographs, the processing
protocol B flags up to $23\%$, and Sb flags
at most $2.7\%$.
%
Among the remaining 755 photographs, Sb flags
at most one image per network.
%
These results suggest improved tolerance to
certain acquisition and resampling differences.
However, the evaluation includes genuine images
only and cannot establish joint attack-detection
performance at a transferable operating threshold.
Furthermore, resolution, compression, and image
content are not independently controlled
in these photographic sources.

\paragraph{Replication and variability.}
Finally, we retrain published comparator methods
using their original training recipes across
five seeds and compare them against their
released checkpoints. These experiments reveal substantial seed
variability and differences between released
and retrained models, particularly for GAN-based
morph families. They measure reproducibility
of the original methods rather than the effect
of our proposed training protocol. Full results, including per-family standard
deviations and additional preprocessing variants,
are reported in the supplementary material.

% \paragraph{Summary.}
% Our experiments establish three main findings.
% First, processing cues alone can perfectly separate
% a specific morph benchmark evaluation without any
% morph supervision.
% Second, matched rendering and manipulation controls
% reveal substantial detector responses that cannot
% be attributed uniquely to identity mixing.
% Third, training with matched controls improves
% two-person versus same-person ranking in controlled
% development comparisons and reduces sensitivity
% to some processing artifacts.

% At the same time, the remaining confounds,
% fixed-threshold limitations, and development-exposed
% benchmark comparisons prevent a claim of universal
% shortcut invariance or reliable identity-specific
% morph detection.
